Congressional redistricting occurs constitutionally every decade following the census, yet research has focused almost exclusively on optimizing single-census solutions. This creates a temporal myopia: while we rigorously evaluate district quality in year zero, we ignore stability as demographics shift over the subsequent decade. Districts optimized for 2020 may fragment communities by 2030, requiring disruptive boundary changes at the next redistricting cycle.

This paper addresses a fundamental yet unexplored question: \textbf{Do different graph partitioning methods produce districts with different temporal stability properties?} Specifically, we compare hierarchical recursive bisection against flat n-way partitioning across two census cycles (2010 and 2020), measuring how district boundaries evolve as demographics change.

Our hypothesis derives from structural differences between methods. Recursive bisection creates nested geographic regions through binary tree decomposition—a state first divides into two super-regions, then each subdivides recursively. This hierarchy may provide temporal scaffolding: top-level splits based on fundamental geography (North vs South, Urban vs Rural) persist even as within-region demographics shift. By contrast, n-way partitioning simultaneously optimizes all $k$ districts without hierarchical structure, potentially causing cascading changes across all boundaries when re-optimized for new census data.

We test this hypothesis on five southern states (Alabama, Georgia, Louisiana, Mississippi, South Carolina) with significant minority populations, running both methods on 2010 and 2020 census data. Four temporal stability metrics quantify district continuity: tract reassignment rates, population disruption, boundary stability, and hierarchical coherence change. This multi-census experimental design—requiring 20 independent redistricting runs—enables the first systematic comparison of partitioning method stability.

Our findings reveal substantial stability differences. Recursive bisection maintains 80\% tract retention across the decade versus 70\% for n-way partitioning (+14 percentage point improvement). This advantage stems from hierarchical structure preservation: recursive bisection's top-level geographic splits remain 94\% intact from 2010 to 2020, while n-way shows no structural continuity. Communities of interest experience 14 percentage points less disruption (22\% vs 36\% counties affected).

These results establish a performance-stability tradeoff. Prior work showed n-way achieves marginally better 2020 minority concentration (+4.3 points), but we demonstrate recursive bisection provides substantially better 2020-2030 stability (+14 points). For jurisdictions with comfortable Voting Rights Act (VRA) margins, stability benefits likely outweigh marginal performance costs.

The implications extend beyond redistricting. Any domain requiring periodic re-partitioning (school districts adjusting to enrollment changes, service territories rebalancing, sales regions evolving with markets) faces similar temporal stability considerations. Our findings suggest hierarchical methods provide continuity advantages that flat optimization methods lack—a principle with broad applicability to dynamic partitioning problems.

\paragraph{Research Program Context} This paper contributes to a broader investigation of recursive bisection for congressional redistricting initiated by foundational work establishing algorithmic objectivity and national-scale feasibility~\cite{della-libera-recursive-bisection}, extended through edge-weighted compactness optimization~\cite{della-libera-edge-weighted} and VRA compliance methods~\cite{della-libera-vra-compliance}. Those studies demonstrate that recursive bisection produces high-quality districts within single census cycles, but a critical practical question remains: does the hierarchical structure provide temporal stability advantages across redistricting cycles? This work answers that question by comparing recursive to n-way partitioning across two census decades (2010-2020), measuring district boundary evolution as demographics shift. Related papers in this program validate algorithmic consistency through slice-based cross-census testing~\cite{della-libera-cross-census-validation} and investigate parameter sensitivity to ensure robust performance across varied configurations.
